
\documentclass{article}
\usepackage{iclr2027_conference,times}
\usepackage[T1]{fontenc}
\usepackage[utf8]{inputenc}
\usepackage{amsmath,amssymb,mathtools}
\usepackage{booktabs,array}
\usepackage{microtype}
\usepackage{algorithm}
\usepackage{algpseudocode}
\usepackage{xcolor}
\usepackage{url}
\usepackage[hidelinks]{hyperref}
\hypersetup{pdftitle={Fact-Addressed Residual Interventions for Selective Behavioral Unlearning},pdfauthor={Anonymous authors}}

\newcommand{\R}{\mathbb{R}}
\newcommand{\ind}{\mathbf{1}}
\newcommand{\relu}[1]{\left[#1\right]_{+}}
\newcommand{\norm}[1]{\left\lVert#1\right\rVert_2}
\newcommand{\bank}{\Delta E}
\newcommand{\rowvec}{\Delta e}
\newcommand{\router}{\mathcal{R}}
\newcommand{\eps}{\varepsilon}
\DeclareMathOperator*{\argmax}{arg\,max}
\DeclareMathOperator*{\argmin}{arg\,min}
\DeclareMathOperator{\NLL}{NLL}
\DeclareMathOperator{\clip}{clip}

% Keep true for the evidence supporting this development-stage abstract.
% Set false when transplanting only the abstract/method into a larger paper.
\newif\ifincludedevelopmentevidence
\includedevelopmentevidencetrue

\title{Fact-Addressed Residual Interventions\\for Selective Behavioral Unlearning}
\author{Anonymous authors}
% Do not enable \iclrfinalcopy for a double-blind submission.

\begin{document}
\maketitle
% Preserve the official review header with recent fancyhdr versions.
\lhead{Under review as a conference paper at ICLR 2027}

\begin{abstract}
Behavioral unlearning of designated factual associations should suppress those
associations while leaving requests about other entities on the base-model
path. We present sparse, fact-addressed residual interventions in a frozen
language model. Each unique protected association receives one trainable
residual vector, while a frozen subject-and-context router selects at most one
vector from the natural request. Positive and negative contextual prototypes
determine whether to intervene; ambiguous requests follow the base-model path.
The selected vector is added at a single transformer block at the original
request boundary, which remains fixed during teacher-forced scoring and
reference generation. Row-wise optimization suppresses sensitive-answer
likelihood without training toward a counterfactual factual answer. In
single-model, development-seed experiments on MCF, zsRE, MQuAKE, and a
probe-assisted RWKU extension, the method uses 153,600--322,560 trainable
parameters. It attains MCF rewrite/paraphrase answer-probability scores of
$0.00017\%$/$0.00059\%$, zero sensitive-token accuracy on zsRE rewrites
and rephrases, and zero direct atomic sensitive-token accuracy on MQuAKE.
An association-level audit reveals that 240 of 1,594 nominal MQuAKE retain
records duplicate protected associations; on the association-disjoint subset,
direct retain accuracy is 0.111 percentage points below the base model.
Measured router activation defines the method's scope: the router abstains on
all 500 MCF neighborhood prompts, but on RWKU it activates on 90\% of held-out
probes about protected people that no trained row represents, and on 4.7\% of
neighbor probes. Intervention is therefore subject-scoped rather than confined
to the trained associations. Same-subject locality and RWKU generated-answer
recovery remain limitations. These results support subject-scoped
retrieval-expression control, not deletion of knowledge from pretrained
weights; confirmatory multi-seed and multi-model evaluation remains
outstanding.
% TODO(kp759): the RWKU activation figures (129/143 held-out, 13/279 neighbor)
% come from outputs/rwku_fact_assoc_router_v2_seed1_direct. Confirm this is the
% run behind the RWKU table, or regenerate the table from it.
\end{abstract}

\section{Methodology}
\label{sec:method}

\subsection{Problem setting and protected association identity}
\label{sec:setting}

Let $p_{\theta_0}$ be a pretrained causal language model. We keep $\theta_0$
fixed, including input embeddings, attention and MLP weights, normalization
parameters, and the language-model head. A protected association is
$a_i=(s_i,r_i,o_i)$, where $s_i$ is a subject, $r_i$ is a symbolic relation or
natural request context, and $o_i$ is the sensitive answer. The goal is to
reduce expression of $o_i$ on requests concerning $a_i$, while preserving
permitted behavior. This is a \emph{behavioral} objective for the model with its
controller attached, not a claim of equivalence to retraining without the
protected information.

The protected set contains $N$ unique associations, not necessarily $N$ raw
records or $N$ distinct subjects. MCF uses sampled factual associations from
Multi-CounterFact. The zsRE adapter uses the direct
request, subject, and sensitive answer in the MEND evaluation setting MQuAKE identities are normalized
$(\text{subject},\text{relation\_id},\text{target\_true})$ triples
duplicate atomic records share a residual. RWKU uses
$(\text{subject},\text{selected query/context},\text{sensitive answer})$.
Its Batch-50 adaptation supplies ten selected probes for each of five people,
and is explicitly \emph{probe-assisted}, unlike native target-only RWKU
 Source-record identifiers establish provenance but are
not supplied to the runtime router.

We allocate a zero-initialized bank
$\bank=[\rowvec_1^\top;\ldots;\rowvec_N^\top]\in\R^{N\times d}$.
Only these rows receive gradients. Frozen contextual prototypes and threshold
metadata are additional stored state, not additional trainable parameters.
In the present Llama-3.2-3B-Instruct experiments, $d=3072$ and $N$ is 50 for
MCF, zsRE, and RWKU, and 105 for MQuAKE. The corresponding residual parameter
counts are $Nd=153{,}600$ and $322{,}560$. No vocabulary expansion, private
fact token, or pretrained-weight modification is introduced.

\subsection{A single intervention at the original request boundary}
\label{sec:boundary}

Let $x=(x_0,\ldots,x_{b-1})$ be the tokenized, formatted request, and let
$h_{\ell,t}$ denote the output of transformer block $\ell$ at position $t$,
before our intervention. Both block and token indices are zero-based. The
frozen configuration uses $\ell=19$. The router returns
$\router(x)\in\{1,\ldots,N\}\cup\{\bot\}$, with $\bot$ denoting no selected
association. Writing $g_i(x)=\ind[\router(x)=i]$, we apply
\begin{equation}
 h'_{\ell,t}
 =h_{\ell,t}+\ind[t=b-1]\sum_{i=1}^{N}g_i(x)\rowvec_i,
 \qquad \sum_{i=1}^{N}g_i(x)\leq 1.
 \label{eq:intervention}
\end{equation}
The remaining blocks and frozen LM head process the altered representation.
Only one position is directly changed at the intervention block; downstream
causal processing can propagate its effect to subsequent positions.

The addressing representation is computed \emph{before} the selected residual
is added:
\begin{equation}
 q(x)=\frac{h_{\ell,b-1}(x)}{\norm{h_{\ell,b-1}(x)}}.
 \label{eq:query}
\end{equation}
Thus the same frozen representation supports addressing throughout row
optimization. The runtime trigger does not require $o_i$, candidate-answer
suffixes, or an externally supplied association index.

For a continuation $y$, we write $p_{\theta_0,\bank}(y\mid x;b)$ to make the
original boundary explicit. Every teacher-forced context uses the same $b$,
even when answer-prefix tokens have been appended. The uncached reference
generator likewise recomputes the growing sequence while intervening at
$b-1$, not at the newest token. Prompt padding is excluded from subject
eligibility and boundary indexing. Tokenization-specific training adapters
reconstruct the scored continuation contexts while checking that the original
request remains the observed prefix. Cached generation is not part of the
validated reference implementation.

\paragraph{Inactive-path preservation.}
For deterministic execution with the same tokenization and boundary, router
abstention makes Equation~\eqref{eq:intervention} the identity. Consequently,
\begin{equation}
 \router(x)=\bot
 \quad\Longrightarrow\quad
 p_{\theta_0,\bank}(y\mid x;b)=p_{\theta_0}(y\mid x)
 \label{eq:inactive}
\end{equation}
for the fixed-boundary continuation. This is a structural property of the
inactive path, not a guarantee that the router rejects every permitted
request. Preservation on routed inputs must be evaluated separately. Removing
the controller exposes the unchanged base model; the method therefore does
not establish irreversible weight-level erasure.

\subsection{Router V2: subject eligibility with contextual confirmation}
\label{sec:router}

\paragraph{Candidate set.}
Let $\mathcal{S}_i$ contain the permitted tokenized subject surfaces for
association $i$. We construct
\begin{equation}
 \mathcal{C}(x)=\{i:\text{a surface in }\mathcal{S}_i
                  \text{ occurs in the valid request prefix }x\}.
 \label{eq:candidates}
\end{equation}
Subject occurrence establishes eligibility only. Unlike Router V1, a unique
eligible subject does not bypass context confirmation. RWKU additionally
permits a surname surface derived from the supplied public name; this is a
lexical eligibility rule, not an entity-resolution guarantee.

\paragraph{Frozen prototypes.}
For each association, positive prompts $\mathcal{X}_i^+$ are drawn only from
its allowed training-visible requests. The MCF adapter includes its canonical
rewrite and authored fitting phrasings; the direct-token adapters use their
registered direct requests. We encode them with Equation~\eqref{eq:query} to
form $P_i$. Negative prompts first use other protected direct contexts for the
same subject. Remaining controls are constructed deterministically by
transplanting the subject into other training-visible prompts. Exact collisions
with an association's positive prompts are excluded. Up to 12 negative prompts
are encoded into $N_i$. These controls are intended to distinguish contexts;
the construction itself does not prove that every synthetic negative is a
semantically valid different relation. Official paraphrases, locality probes,
retain records, and utility or adversarial evaluation data do not construct
these prototype banks.

For normalized query $q$, define
\begin{align}
 u_i(q)&=\max_{p\in P_i}\cos(q,p),
 &v_i(q)&=\max_{n\in N_i}\cos(q,n),\nonumber\\
 d_i(q)&=u_i(q)-v_i(q).
 \label{eq:score}
\end{align}
The implementation stores an absolute threshold $\alpha_i=-1$, which disables
absolute-similarity rejection. Qualification therefore depends on the relative
context score $d_i$, not on a separately tuned positive-cosine floor.

\paragraph{Threshold fitting.}
Let $d^+_{i,\min}$ be the minimum margin on the positive fitting prompts and
$d^-_{i,\max}$ the maximum on the negative fitting prompts. The frozen rule is
\begin{equation}
 \tau_i=\clip_{[-2,2]}\begin{cases}
 d^-_{i,\max}+\lambda\bigl(d^+_{i,\min}-d^-_{i,\max}\bigr),
       &d^-_{i,\max}<d^+_{i,\min},\\
 d^+_{i,\min}-\sigma,&\text{otherwise},
 \end{cases}
 \quad \lambda=0.10,\quad\sigma=0.02.
 \label{eq:threshold}
\end{equation}
The separable-case threshold lies above every observed fitting negative and
below every observed fitting positive. In the nonseparable case the fallback
prioritizes positive fitting coverage and can admit negatives. These are
\emph{in-sample} properties: the current builder reuses prototype-bearing
prompts for threshold fitting, rather than estimating independent calibration
error. The 10\% operating point was selected during seed-1 development.
Confirmatory experiments freeze this rule and its hyperparameters; numerical
thresholds are constructed from each new bank's allowed fitting data without
held-out retuning.

\paragraph{Selection and abstention.}
Let $\mathcal{Q}(x)=\{i\in\mathcal{C}(x):d_i(q(x))\geq\tau_i\}$.
If $\mathcal{Q}(x)$ is empty, the router returns $\bot$. Otherwise, let $i_1$
and $i_2$ be the highest- and second-highest-scoring qualifying candidates.
The selected route is
\begin{equation}
 \router(x)=
 \begin{cases}
 i_1,&|\mathcal{Q}(x)|=1,\\
 i_1,&|\mathcal{Q}(x)|>1\ \text{and}\ d_{i_1}-d_{i_2}\geq\mu,\\
 \bot,&\text{otherwise},
 \end{cases}
 \qquad\mu=0.02.
 \label{eq:selection}
\end{equation}
The separation is computed among \emph{qualifying} candidates. An ambiguous
request receives no residual; this is router abstention, not necessarily a
verbal refusal by the base model. At most one association is modified per
request, so simultaneous control of multiple forbidden facts is not guaranteed.

\begin{algorithm}[t]
\caption{Fixed-boundary inference with Router V2}
\label{alg:inference}
\begin{algorithmic}[1]
\Require Frozen model $\theta_0$, bank $\bank$, frozen prototypes and thresholds,
formatted request $x$, maximum continuation length $T$
\State $b\gets |x|$; $z\gets x$
\For{$k=1,\ldots,T$}
    \State Run frozen blocks through block $\ell$ on $z$ without a KV cache
    \State Read $q\gets\operatorname{normalize}(h_{\ell,b-1})$ before intervention
    \State Compute $\mathcal{C}(x)$ from request tokens only
    \State Select $i\gets\router(x)$ using Equations~\eqref{eq:score}--\eqref{eq:selection}
    \If{$i\neq\bot$}
        \State $h_{\ell,b-1}\gets h_{\ell,b-1}+\rowvec_i$
    \EndIf
    \State Run the remaining frozen blocks and LM head
    \State Append the greedy next token to $z$
    \If{the configured stopping token is generated}
        \State \textbf{break}
    \EndIf
\EndFor
\State \Return the continuation $z_{b:}$
\end{algorithmic}
\end{algorithm}

\subsection{Row-wise sensitive-answer suppression}
\label{sec:optimization}

\paragraph{Direct-token objective.}
For zsRE, MQuAKE, and the RWKU direct stage, let
$\mathcal{T}_i=\{(c_{ijt},b_{ij},y_{ijt})\}$ contain the allowed teacher-forced
contexts for association $i$. The context $c_{ijt}$ is reconstructed by the
dataset adapter for answer token $t$ of fitting request $j$, while $b_{ij}$
remains the original request boundary. Define
\begin{align}
 \ell_{ijt}(\rowvec_i)
 &=-\log p_{\theta_0,\bank}(y_{ijt}\mid c_{ijt};b_{ij}),\nonumber\\
 Q_i^{\mathrm{tok}}(\rowvec_i)
 &=\max_{j,t}\exp(-\ell_{ijt}(\rowvec_i)),\nonumber\\
 L_i^{\mathrm{tok}}
 &=\relu{A-\min_{j,t}\ell_{ijt}(\rowvec_i)},
 \qquad A=-\log\eps,\quad\eps=10^{-6}.
 \label{eq:tokenloss}
\end{align}
The gradient targets the currently most probable sensitive token. This is
stronger than requiring only one answer token to cease being top-1. Feasibility
is checked using the strict condition $Q_i^{\mathrm{tok}}<\eps$; a zero hinge
at equality is not reported as satisfying that strict condition. Feasible rows
are skipped in these direct-only trainers. No counterfactual answer or
abstention completion is optimized in those three direct stages.

\paragraph{MCF objective and abstention phase.}
MCF retains its existing worst-phrasing, mean-token-NLL objective. For a fitting
view $j$ with $m_{ij}$ scored answer tokens,
\begin{equation}
 \bar\ell_{ij}=\frac{1}{m_{ij}}\sum_t\ell_{ijt},\qquad
 Q_i^{\mathrm{geo}}=\max_j e^{-\bar\ell_{ij}},\qquad
 L_i^{\mathrm{geo}}=\relu{A-\min_j\bar\ell_{ij}}.
 \label{eq:mcfloss}
\end{equation}
Here $e^{-\bar\ell}$ is a geometric-mean token probability, not the largest
constituent-token probability or the full-answer probability. Before a row
satisfies $Q_i^{\mathrm{geo}}<\eps$, proposals follow the suppression loss.
Afterward, the phase-lexicographic trainer minimizes the mean NLL of the
ordinary completion ``I don't know.'' while maintaining the suppression
constraint. The abstention weight is 1.0. This secondary completion supplies an abstention preference rather than
counterfactual factual supervision. The direct-token adapters omit this phase.

\paragraph{Proposals and acceptance.}
Each row has a separate Adam optimizer with learning
rate 0.05. Only the selected row is updated at a row step, and its gradient
norm is clipped to 1.0. Let $\xi_i$ be the displacement proposed by Adam and
$Q_i$ the adapter's current worst-case probability score. We rescale
\begin{equation}
 \tilde\xi_i=\xi_i\min\left\{1,\frac{r(Q_i)}{\max(\norm{\xi_i},10^{-30})}\right\},
 \qquad
 \rowvec_i^{(k)}=\rowvec_i+2^{-k}\tilde\xi_i,\quad k=0,\ldots,12.
 \label{eq:backtrack}
\end{equation}
The radius is 1.00, 0.35, 0.08, or 0.02 as the current probability decreases
through the $10^{-3}$, $10^{-5}$, and $10^{-6}$ breakpoints; exact boundary
conventions are specified in Appendix~\ref{app:settings}. This is a bound on
the \emph{step}, not on the accumulated residual norm.

Candidates are rescored through the complete edited model. The direct-token
trainers accept a strict decrease in $Q_i^{\mathrm{tok}}$ and select the best
acceptable backtracking candidate. MCF uses lexicographic candidate ordering:
$(Q_i^{\mathrm{geo}},L_i^{\mathrm{abst}})$ before feasibility and
$(L_i^{\mathrm{abst}},Q_i^{\mathrm{geo}})$ afterward, disallowing violation of a
locked suppression constraint. If no candidate is accepted, both the row and
its optimizer state are restored. These checks constrain sensitive-answer
behavior on finite fitting contexts; they impose neither a retained-output KL
budget nor a universal locality guarantee.

Rows are visited in seeded shuffled sweeps. The nominal budget is 30 attempted
updates per unique association, subject to the recorded wall-time and stall
limits. Checkpoint evaluations occur at complete sweeps. The direct-token
trainers select the smallest maximum training-token probability. MCF selects
using its train/development worst-view feasibility and secondary abstention
score; authored development prompts supply no gradients or prototype fitting,
but do influence checkpoint selection. All trainers restore the selected
checkpoint. Budget termination, zero measured accuracy, and satisfaction of
$\eps$ are reported as distinct outcomes.

\subsection{Evaluation contract and association-aware locality}
\label{sec:evaluation}

The current evidence consists of a single development seed on one pretrained
model, with separately trained banks for four benchmark adapters. It does not
constitute transfer of one trained bank across datasets or model families.
Training uses float32; final evaluation uses the recorded deployment precision,
with base and edited models matched in tokenizer and scoring protocol. The
completed snapshot does not include Router V2 rollout hardening or confirmatory
multi-model runs.

\paragraph{Noninterchangeable forgetting metrics.}
For MCF, complete-answer probability is
\begin{equation}
 P_{\theta_0,\bank}(o_i\mid x;b)
 =\exp\!\left(-\sum_{t=1}^{m_i}\ell_{it}\right).
 \label{eq:sequenceprob}
\end{equation}
The reported Eff and Gen are 100 times its mean on canonical and paraphrased
requests, respectively, averaging phrasings within each case before averaging
cases. EOS is not added to the scored answer. Neighborhood Spe is all-answer-token
teacher-forced top-1 accuracy. We retain token-geometric likelihood,
all-token correctness, and true-versus-counterfactual preference as separate
diagnostics. A display-zero check requires unrounded Eff and Gen below
$0.005\%$ and zero all-token sensitive-answer accuracy; it does not claim
mathematically zero probability.

The zsRE adapter instead reports case-macro sensitive-token top-1 accuracy
for rewrites and rephrases. MQuAKE reports direct atomic sensitive-token
accuracy and a separate held-out atomic-question score, AtomicGen; neither is
a multi-hop guarantee. RWKU reports generated-answer recovery and ROUGE-L,
alongside teacher-forced sensitive-token accuracy, for same-50, held-out,
paraphrase, adversarial, and neighbor groups. These scores cannot be pooled as
though Eff or Gen had one definition across all four datasets. The RWKU
snapshot is direct-stage only, and its zero sensitive-token accuracy coexists
with nonzero generated-answer recovery.

\paragraph{Overlap-stratified retain evaluation.}
Let $\mathcal{F}$ be the set of normalized protected atomic identities and
$\mathcal{R}_{\mathrm{raw}}$ the nominal retain records. We identify
\begin{equation}
 \mathcal{R}_{\cap}
 =\{r\in\mathcal{R}_{\mathrm{raw}}:a(r)\in\mathcal{F}\},\qquad
 \mathcal{R}_{\setminus}
 =\mathcal{R}_{\mathrm{raw}}\setminus\mathcal{R}_{\cap}.
 \label{eq:overlap}
\end{equation}
MQuAKE source-instance disjointness does not establish atomic-association
disjointness. We therefore report the unchanged overall benchmark retain
score \emph{and} an association-disjoint diagnostic. Within
$\mathcal{R}_{\setminus}$, we separate same-subject/different-relation records
from subject-disjoint records. Metadata-level subject disjointness does not
exclude lexical overlap in a prompt. Exact protected overlaps are reported,
not silently removed from the official score; the stratification is an
analysis of what that score measures, not a replacement benchmark.

The development audit finds 240 exact protected overlaps among 1,594 nominal
retain atomic records, leaving 1,354 association-disjoint records. Of these,
81 share a protected subject but have another relation, and 1,273 are
subject-disjoint. This decomposition is especially important when interpreting
activation rates: intervention on an exact protected association is consistent
with the specified suppression target, whereas permitted same-subject
relations provide a harder locality test. The overlap audit was performed
after observing development results and must be labeled accordingly.

\paragraph{Utility and route diagnostics.}
A conventional whole-sequence PPL call can miss a last-position-only edit
because its final logits do not predict an observed token. We therefore also
score each target from its observed prefix, with the boundary at the position
predicting that target, normalizing by the actual number of scored targets.
This is explicitly a \emph{prefix-reset next-token} measurement, not the same
boundary schedule as a multi-token response to one fixed request. The present
99-target WikiData fixture has zero route activity; equal PPL there supports
inactive-path preservation only. For generated answers we retain the
fixed-boundary schedule of Algorithm~\ref{alg:inference}. Route activation,
correct association selection, suppression, and retained behavior are distinct
measurements; none substitutes for the others.

\paragraph{Measured activation scope.}
Route activation is reported directly rather than inferred from unchanged
utility scores, because an identical Spe or PPL across methods is equally
consistent with a local residual and with a router that never fires. On MCF,
Router~V2 activates on 0 of 500 neighborhood and 0 of 1,000 retain prompts
(95\% Wilson upper bounds 0.76\% and 0.38\%) and selects the correct
association on 147 of 150 rewrite and paraphrase prompts, so its neighborhood
Spe is the inactive path of Equation~\eqref{eq:inactive} by measurement rather
than by assumption. RWKU differs. Its held-out Level-1, Level-2 and paraphrase
probes concern the protected people but are content-disjoint from the trained
probes, so no residual row represents them. Router~V2 nonetheless activates on
129 of 143 of these probes (90.2\%, 95\% CI 84.2--94.1\%), in each case routing
to a trained row of the same person, and on 13 of 279 neighbor probes (4.7\%,
95\% CI 2.7--7.8\%). Subject eligibility followed by context confirmation
therefore scopes intervention to the protected subject rather than to the
individual trained association. A decomposition that supplies the ground-truth
row separates routing from actuation on the trained associations: on the
RWKU same-50 probes, Router~V2 and ground-truth routing yield identical
generated-answer recovery, so the residual recovery there is an actuation
effect.
% TODO(kp759): add MQuAKE activation on the 81 same-subject/different-relation
% retain records; it is the same question on a second benchmark.

\subsection{Claim scope and frozen configuration}
\label{sec:claims}

The method separates natural-input addressing from a small trainable actuator.
It neither locates exclusive storage of a fact nor removes information from
$\theta_0$. Literal subject eligibility limits alias and coreference coverage;
prototype controls do not certify semantic scope; ambiguous or unmatched
forbidden requests can fall back to a knowledgeable base model. Single-route
selection also does not establish selective handling of mixed multi-fact
requests. In particular, the router must not be read as intervening only on
the associations represented by trained rows: on RWKU it activates on most
held-out requests about a protected person
(Section~\ref{sec:evaluation}). Whether suppression of those unlisted facts
reflects subject-conditioned transfer or generic disruption of generation is
tested separately, by giving the same held-out probes an equal number of
residual rows trained on other people's facts.
% TODO(kp759): insert the cross-person best-of-K control result here. Until it
% reports, do not describe held-out RWKU suppression as transfer.
The current results must therefore be interpreted as subject-scoped
retrieval-expression control under the stated interface.

The development snapshot fixes block 19, the contextual threshold rule,
negative-control count, ambiguity margin, optimizer, and step-radius schedule
before confirmatory evaluation. MCF/zsRE/MQuAKE sampling seeds and RWKU's
deterministic five-person windows are different replication units and should
be reported separately. The numerical development evidence supporting the
abstract\ifincludedevelopmentevidence
 is retained in Appendix~\ref{app:evidence} and\fi
 does not stand in for the planned seeds 1--10 or additional model families.


\appendix
\section{Implementation settings}
\label{app:settings}

All entries below describe the frozen development configuration, not a
hyperparameter search result or an optimality claim. The residual bank starts
at zero. Contextual prototypes are frozen after construction, and the
absolute-similarity floor is disabled. With $\alpha=-1$, only the relative context
score and ambiguity separation reject subject-eligible requests.

\begin{table}[ht]
\centering
\caption{Frozen row-training budgets. An update denotes a row-step opportunity,
not a full-bank update; feasible rows may be skipped.}
\label{tab:budgets}
\small
\begin{tabular}{@{}lrrrrr@{}}
\toprule
Adapter & Raw records & Rows & Nominal steps & Time cap (s) & Max length\\
\midrule
MCF & 50 & 50 & 1,500 & 3,600 & 512\\
zsRE & 50 & 50 & 1,500 & 3,600 & 512\\
MQuAKE & 127 atomic & 105 & 3,150 & 7,200 & 512\\
RWKU & 50 probes & 50 & 1,500 & 7,200 & 4,096\\
\bottomrule
\end{tabular}
\end{table}

The optimizer uses Adam at learning rate 0.05, gradient clipping at 1.0,
12 halvings of the initial proposal, and a stall limit of 150 rejected row
steps. The radius helper checks breakpoints with strict greater-than
comparisons. Its executable convention is
\begin{equation}
 r(Q)=\begin{cases}
 1.00,&Q>10^{-3},\\
 0.35,&10^{-5}<Q\leq10^{-3},\\
 0.08,&10^{-6}<Q\leq10^{-5},\\
 0.02,&Q\leq10^{-6}.
 \end{cases}
 \label{eq:radius}
\end{equation}
% Author review: the snapshot table writes >= at the upper breakpoints;
% The radius equation follows radius_for_probability() in the pinned source.
Router V2 uses up to 12 negative controls, $\lambda=0.10$,
$\sigma=0.02$, and $\mu=0.02$. MCF additionally retains a 0.90 authored-development
route-recall preflight floor and two post-feasibility checkpoint
checks. MCF's legacy gate-slack configuration field of 0.04 is not the
context-margin slack $\sigma$.

\ifincludedevelopmentevidence
\section{Development evidence supporting the abstract}
\label{app:evidence}

These are reported seed-1 results from the frozen research snapshot, not
independently reproduced measurements. Router V2 was selected during this
development cycle. No multi-seed uncertainty, multi-model result, or superiority
over external unlearning baselines is claimed here. 






\begin{table}[t]
\centering
\small
\caption{MCF seed-1 development results on Llama-3.2-3B-Instruct.
Eff and Gen are complete sensitive-answer probabilities (percent) on rewrite
and paraphrase requests. Spe is neighborhood all-answer-token teacher-forced
accuracy (percent). PPL is the prefix-reset next-token measurement.}
\label{tab:mcf-seed1-eff-gen-spe-ppl}
\begin{tabular}{@{}lrrrr@{}}
\toprule
Method & Eff $\downarrow$ & Gen $\downarrow$ & Spe $\uparrow$ & PPL $\downarrow$ \\
\midrule
Frozen base & 12.266065 & 7.766598 & 20.400000 & 11.312981 \\
Router V1   & 0.000066955 & 0.000230665 & 20.400000 & 11.312981 \\
Router V2   & 0.000178522 & 0.000595133 & 20.400000 & 11.312981 \\
\bottomrule
\end{tabular}
\end{table}

\begin{table}[t]
\centering
\small
\caption{zsRE seed-1 development results on Llama-3.2-3B-Instruct.
Eff and Gen are sensitive-answer-token teacher-forced top-1 accuracy (percent)
on direct and rephrased requests. Spe is locality-token accuracy (percent).
}
\label{tab:zsre-seed1-eff-gen-spe-ppl}
\begin{tabular}{@{}lrrrr@{}}
\toprule
Method & Eff $\downarrow$ & Gen $\downarrow$ & Spe $\uparrow$ & PPL $\downarrow$ \\
\midrule
Frozen base & 29.966667 & 27.900000 & 32.323179 & 11.312981 \\
Router V1   & 0.000000 & 0.000000 & 32.323179 & 11.312981 \\
Router V2   & 0.000000 & 0.000000 & 32.323179 & 11.312981 \\
\bottomrule
\end{tabular}
\end{table}

\begin{table}[t]
\centering
\small
\caption{MQuAKE seed-1 development results on Llama-3.2-3B-Instruct.
Eff is direct atomic sensitive-token accuracy (percent).
The Gen column denotes the reported AtomicGen extension: held-out
atomic-question sensitive-token accuracy (percent), not multi-hop accuracy.
No separate Spe score is defined in the supplied result contract; a dash means
not reported, not zero. Overall and association-disjoint retain scores remain
separate analyses and must not be substituted for Spe.}
\label{tab:mquake-seed1-eff-gen-spe-ppl}
\begin{tabular}{@{}lrrrr@{}}
\toprule
Method & Eff $\downarrow$ & Gen$^{\dagger}$ $\downarrow$ & Spe & PPL $\downarrow$ \\
\midrule
Frozen base & 73.063367 & 50.189351 & --- & 11.312981 \\
Router V1   & 0.000000 & 1.771654 & --- & 11.312981 \\
Router V2   & 0.000000 & 2.952756 & --- & 11.312981 \\
\bottomrule
\end{tabular}
\par\smallskip
\begin{minipage}{0.98\linewidth}
\footnotesize $^{\dagger}$Gen = AtomicGen; this is not a native multi-hop result.
\end{minipage}
\end{table}

\begin{table}[t]
\centering
\small
\caption{RWKU probe-assisted Batch-50 seed-1 development results on
Llama-3.2-3B-Instruct, direct stage only. To provide the requested column layout,
Eff denotes same-50 generated-answer recovery; Gen denotes held-out Level-2
paraphrase recovery; Spe denotes neighbor recovery. All three are percentages.
These are explicitly mapped RWKU measures, not the MCF metric definitions.
No rollout-hardened checkpoint is included.}
\label{tab:rwku-seed1-eff-gen-spe-ppl}
\begin{tabular}{@{}lrrrr@{}}
\toprule
Method & Eff$^{\ddagger}$ $\downarrow$ & Gen$^{\ddagger}$ $\downarrow$ & Spe$^{\ddagger}$ $\uparrow$ & PPL $\downarrow$ \\
\midrule
Frozen base & 60.000000 & 62.000000 & 69.175630 & 11.312981 \\
Router V1   & 10.000000 & 44.000000 & 65.591398 & 11.312981 \\
Router V2   & 10.000000 & 44.000000 & 65.591398 & 11.312981 \\
\bottomrule
\end{tabular}
\par\smallskip
\begin{minipage}{0.98\linewidth}
\footnotesize $^{\ddagger}$Eff = same-50 recovery; Gen = held-out Level-2
paraphrase recovery; Spe = neighbor recovery. The other held-out levels and
Level-3 adversarial recovery are not replaced by this three-column summary.
\end{minipage}
\end{table}


\end{document}
